\documentclass[a6paper, 10pt, twoside]{article}
%\usepackage[T1]{fontenc}
\usepackage[british]{babel}
\usepackage[utf8]{inputenc}
\usepackage{float, graphicx,amsmath,amsfonts,cite,enumerate,tabularx}
\usepackage[final]{pdfpages}
\usepackage{wrapfig}
\usepackage[margin=0.3in]{geometry}
\usepackage{sidspaltHack}
\usepackage{digital}
\usepackage{ulem}


\setlength{\evensidemargin}{-0.47in}
\setlength{\oddsidemargin}{-0.37in}
\setlength{\textwidth}{215pt}

\pagestyle{empty}

\begin{document}
\nysida{16}{}
\noindent
\chaptertitlenobr{$\Phi\phi$}{Visor från spex}
\small
\vspace{10pt}

\noindent
\small
I denna års upplaga av sångboken inför vi ett nytt kapitel: Phi! Här samlar vi de bästa låtarna från den senaste uppsättningen av Fysikalen. 
\auth{Christoffer Ousbäck \\Sångboksansvarig 2024}

\nysida{16}{1}

\begin{center}
    \songtitle{$\phi1$}{Irwins Zoolo}
    \mel{Zwampen}
\end{center}
\begin{lyrics}
För den som har rätt kunskap\\
Och kan zoologi\\
Finns det alltid mat i\\
Naturens skafferi
\vspace{5pt}\\
En termit dränkt i sprit retar din aptit (Sweet!)\\
Ta en ters, sen en bärs till koala-färs (Äsch)\\
Dingo-biff sätter piff på din aperitif (Sniff)\\
En vombat på ett fat, det är delikat (Mat)
\vspace{5pt}\\
När magen börjar kurra\\
Och hungern knackar på\\
Då kan man koka burra\\
Eller äta rå
\vspace{5pt}\\
Bäst ragu, det gör du av en känguru (Bu!)\\
Vilda bin, myrpiggsvin, äts med spetsat vin (Aj)\\
Vallaby i fondue gör dig frisk och kry (Fy)\\
Efterrätt blir en lätt näbbdjursomelett (Fett)
\vspace{5pt}\\
Man tager vad man haver\\
Det sa Cajsa Warg\\
Så jag tog vad hon hade\\
Då blev hon ganska arg
\vspace{5pt}\\
Spindelben i gelén piffar upp buffén\\
Pungvargsspett gör dig mätt, stekt i brunormsfett (Hett)\\
Paddbuljong äts med tång, max en enda gång (Gång)\\
Pytontarm har sin charm om den ätes varm\\
Vilken tur, vår natur har så goda djur (Tur)
\end{lyrics}

\nysida{16}{2}
\begin{center}
    \songtitle{$\phi2$}{Hunt hade en gång}
    \mel{Jag hade en gång en båt}
\end{center}
\begin{lyrics}
Jag hade en stridskamrat\\
Emuerna åt som mat\\
Men det var för länge sen, så länge sen\\
Var är han nu? (Var är han nu?)\\
I en emu! (I en emu!)\\
Han åts som förrätt, av en emu
\vspace{5pt}\\
Jag hade en skyttegrupp\\
Men dem åt en emu upp\\
Ja det var så länge sen, så länge sen\\
Var är de nu? (Var är de nu?)\\
I en emu! (I en emu!)\\
De fyllde magen, på en emu
\vspace{5pt}\\
Jag hade en bataljon\\
Som blev en emu-ranson\\
Och det var en vecka sen, de åt upp den\\
Var är de nu? (Var är de nu?)\\
Ser ni emun? (Ser ni emun?)\\
Så börja springa bort från emun!
\end{lyrics}

\nysida{16}{3}
\begin{center}
    \songtitle{$\phi3$}{The Phantoms of the Outback}
    \mel{Phantom of the Opera}
\end{center}
\begin{lyrics}
\textit{Hunt:}\\
Jag hör en fågelsång\\
Jag hört förut\\
Ett anfall är på gång\\
Det är akut\\
Är det så här det känns\\
Med déjà vu?\\
Emuerna från skuggan träder fram\\
För sent att fly
\vspace{5pt}\\
\textit{Amelia:}\\
Finns ingen mening med\\
Att kämpa mot\\
Och ingen chans för fred\\
Mot detta hot\\
Att följa med er hit,\\
Det ångrar jag\\
Emuerna från skuggan skrider fram\\
För sent att dra
\vspace{5pt}\\
\textit{Kelly:}\\
Får aldrig veta vad\\
Som hänt mitt plan\\
\textit{Irwin:}\\
Min plan för jakt gick ej\\
Som jag är van\\
\textit{Kelly + Irwin:}\\
Nu jagas vi av dem\\
Säg, vilken pärs\\
Emuerna från skuggan strögar fram\\
Vi tar en bärs
\vspace{5pt}\\
(Se upp)\\
(För ruggiga emuer)\\
(Ge upp)\\
(För skuggiga emuer)
\vspace{5pt}\\
\textit{George:}\\
Jag har en ny taktik\\
Ett mästerdrag\\
Du distraherar dem\\
Så rymmer jag\\
En societetsperson\\
Är nog mer mör\\
Sjung "We shall overcome" så går det bra\\
Annars, ni dör
\vspace{5pt}\\
\textit{Dorothy:}\\
Att va i vildmarken\\
Är ej min grej\\
Men för att hitta dig\\
Jag offrar mig\\
Mitt eget öde, det\\
Är sekundärt\\
Jag fick en liten chans att hitta dig\\
Det var det värt
\vspace{5pt}\\
(Se upp)\\
(För hungriga emuer)\\
(Ge upp)\\
(Står chanslös mot emuer)
\end{lyrics}

\nysida{16}{4}
\begin{center}
    \songtitle{$\phi4$}{Starta en chark-business}
    \mel{Bara bada bastu}
\end{center}
\begin{lyrics}
Nå jaa
\vspace{5pt}\\
Vi har en affärsidé\\
Som vi tror kan bli stor succé\\
Åk till landet med gran och en\\
Bygg en gård där på Österlen
\vspace{5pt}\\
(Oh, eh-oh, eh-oh)\\
Mata djuret med hö och halm\\
(Oh, eh-oh, eh-oh)\\
Sen sälj det skuret på Östermalm\\
(Yxa, slafsa, dolme, fauna)
\vspace{5pt}\\
Öppna upp en praktisk charkdisk\\
Säljes av ett barn från dörr till dörr\\
Överklassen ger en peng till klassen när de vill köpa körv\\
Fylla på i Schulmans kylvagn\\
Platea upp med vindruvor och brie
Oh-oh, oh-oh, oh-oh\\
Föreslaget vin:\\
Ca(u)va
\end{lyrics}

\end{document}